% =====================================================================
%  Study A — Spatial Preservation vs. Parameter Count
%  Controlled architecture ablation (results report).
%  Compile with Tectonic ( tectonic StudyA_Report.tex ) or pdfLaTeX.
% =====================================================================
\documentclass[11pt]{article}
\usepackage[a4paper,margin=1in]{geometry}
\usepackage{booktabs}
\usepackage{array}
\usepackage{amsmath,amssymb}
\usepackage{xcolor}
\usepackage{enumitem}
\usepackage[hidelinks]{hyperref}
\usepackage{parskip}
\setlength{\emergencystretch}{3em}

\newcommand{\best}[1]{\textbf{#1}}
\title{\textbf{Study A — Spatial Preservation vs.\ Parameter Count}\\[2pt]
\large A controlled architecture ablation for WiFi indoor localization}
\author{DLoc lightweight localization --- E-JUST graduation project}
\date{\today}

\begin{document}
\maketitle
\vspace{-1.1em}\hrule\vspace{1em}

\textbf{Bottom line.} Across a controlled ladder of nine architectures on the WILD
Jacobs-Hall in-distribution split, \emph{adding parameters does not improve accuracy, but
adding spatial structure does}. Pouring capacity into the decoder of the plain
MobileNetV2 leaves the error essentially flat ($0.881\!\to\!0.883$\,m as parameters rise
$2.02\!\to\!2.92$\,M), whereas adding U-Net skip connections lowers it at \emph{every}
capacity, and a skip model with \emph{fewer} parameters (2.41\,M) beats a no-skip model
with \emph{more} (2.92\,M). Multi-scale features (ASPP) help further (best mean 0.791\,m).
The skip vs.\ no-skip difference is highly significant (paired Wilcoxon $p=1.8\times10^{-11}$).
\textbf{This directly supports the thesis} and hardens the Mamba / 2-D spatial argument.

% =====================================================================
\section{Objective}
% =====================================================================
This study answers the question posed by the supervisor: isolate the contribution of each
architectural ingredient on top of a MobileNetV2 backbone ---
$+\text{decoder}\to+\text{skip connections}\to+\text{multi-scale features}\to
+\text{bottleneck widths}\to+\text{decoder capacities}$ --- and test whether
\emph{spatial preservation contributes independently of parameter count}. A positive result
strengthens the central argument of the project: that \emph{preserving 2-D spatial structure},
not raw model size, is what makes a compressed CSI localizer accurate (the same reason the
1-D Mamba variant, which flattens the 2-D heatmap, fails).

% =====================================================================
\section{Experimental design}
% =====================================================================
A single \emph{configurable} student (\texttt{student\_configurable.py}) exposes four knobs ---
\texttt{use\_skips}, \texttt{decoder\_width}, \texttt{width\_mult} (encoder/bottleneck width),
and \texttt{bottleneck} (\texttt{none}/\texttt{aspp}) --- and reproduces the shipped networks
exactly (no-skip, width 1.0 $=$ the 2.27\,M MobileNetV2; skip, width 1.0 $=$ the 2.41\,M
MobileNetV2-UNet). Thirteen training runs cover nine configurations.

\begin{itemize}[nosep,leftmargin=1.4em]
\item \textbf{Capacity axis $=$ decoder width} (0.5/1.0/2.0). The encoder stays
      ImageNet-pretrained at width 1.0, so this axis varies capacity \emph{without} a
      pretraining confound. Encoder-width (\texttt{width\_mult} 0.5/1.4) is a secondary,
      \emph{from-scratch} axis (no pretraining at non-unit width) and is reported with that caveat.
\item \textbf{Data:} WILD Jacobs-Hall in-distribution split (Fig.\ 10b): 17{,}574 training /
      4{,}260 test samples. Input $4\times161\times361$ AoA--ToF heatmaps; metric is the
      Euclidean peak-to-peak distance $\times0.1$\,m.
\item \textbf{Training:} Adam, lr $10^{-4}$, batch 32, 120 epochs, cosine annealing; MSE loss.
      Best checkpoint selected by \emph{mean} error (the primary metric; the median is
      pixel-quantized to $0.1\sqrt{n}$ and cannot resolve sub-decimetre differences).
\item \textbf{Seeds:} the two headline configurations (no-skip and skip at decoder-width 1.0)
      use 3 seeds (42/123/777) for error bars; the remaining configurations are single-seed
      (trend only). Statistical test: paired Wilcoxon on per-sample errors (seed 42, $n=4260$).
\end{itemize}

% =====================================================================
\section{Results}
% =====================================================================
\begin{table}[h]
\centering
\caption{All nine configurations, ordered by parameter count. Mean error is the primary
metric ($\pm$ sample std over 3 seeds for the two anchors). ``pre'' $=$ ImageNet-pretrained
encoder. Best mean in bold.}
\setlength{\tabcolsep}{4.5pt}\small
\begin{tabular}{lccccccc}
\toprule
Config & Skips & Multi-scale & Pre & Params & FLOPs & Mean (m) & P90 (m) \\
\midrule
skip, enc-w 0.5$^\dagger$        & yes & -- & no  & 0.83M & 0.76G & 1.037           & 1.868 \\
no-skip, dec-w 0.5               & no  & -- & yes & 2.02M & 0.92G & 0.881           & 1.565 \\
skip, dec-w 0.5                  & yes & -- & yes & 2.09M & 1.00G & 0.851           & 1.513 \\
\textbf{no-skip, dec-w 1.0}      & no  & -- & yes & 2.27M & 1.28G & $0.882\pm0.018$ & 1.570 \\
\textbf{skip, dec-w 1.0}         & yes & -- & yes & 2.41M & 1.45G & $0.830\pm0.003$ & 1.435 \\
no-skip, dec-w 2.0               & no  & -- & yes & 2.92M & 2.51G & 0.883           & 1.562 \\
skip, dec-w 2.0                  & yes & -- & yes & 3.19M & 2.85G & 0.866           & 1.513 \\
skip, dec-w 1.0, ASPP            & yes & yes& yes & 3.28M & n/a$^\ddagger$ & \best{0.791} & \best{1.432} \\
skip, enc-w 1.4$^\dagger$        & yes & -- & no  & 4.31M & 2.28G & 0.855           & 1.581 \\
\bottomrule
\end{tabular}
\end{table}
{\footnotesize $^\dagger$ Encoder-width $\neq1.0$: trained from scratch (no ImageNet
pretraining) --- confounded, secondary axis. $^\ddagger$ The FLOP counter does not support the
ASPP module; value omitted.}

\begin{table}[h]
\centering
\caption{Mean error vs.\ capacity, read down each column (pretrained, decoder-width axis).
The no-skip row is \emph{flat} as parameters grow; the skip row is lower at every width.}
\begin{tabular}{lccc}
\toprule
Decoder width & 0.5 & 1.0 & 2.0 \\
\midrule
no-skip \; (params $\to$)  & 2.02M & 2.27M & 2.92M \\
\quad mean error (m)       & 0.881 & 0.882 & 0.883 \;\;\emph{(flat)} \\
\midrule
skip \; (params $\to$)     & 2.09M & 2.41M & 3.19M \\
\quad mean error (m)       & 0.851 & \best{0.830} & 0.866 \;\;\emph{(lower everywhere)} \\
\bottomrule
\end{tabular}
\end{table}

% =====================================================================
\section{Findings}
% =====================================================================
\begin{enumerate}[leftmargin=1.5em,itemsep=4pt]
\item \textbf{Skip connections help, significantly.} At matched decoder-width 1.0, skips lower
      the mean from $0.882\pm0.018$ to $0.830\pm0.003$\,m ($-6.0\%$); the per-sample paired
      Wilcoxon test gives $p=1.8\times10^{-11}$ ($n=4260$). The tiny std of the skip anchor
      ($\pm0.003$) shows the effect is not seed noise.

\item \textbf{Parameter count alone buys nothing (the key result).} The no-skip mean is
      \emph{flat} --- $0.881/0.882/0.883$\,m --- as the decoder grows $2.02\!\to\!2.27\!\to\!2.92$\,M
      parameters (and FLOPs nearly triple, $0.92\!\to\!2.51$\,G). Tripling the decoder's
      capacity does not move the accuracy at all.

\item \textbf{Spatial preservation beats extra parameters directly.} A skip model with
      \emph{fewer} parameters outperforms a no-skip model with \emph{more}: skip dec-w 0.5
      (2.09\,M, 0.851\,m) $<$ no-skip dec-w 1.0 (2.27\,M, 0.882\,m), and skip dec-w 1.0
      (2.41\,M, 0.830\,m) $<$ no-skip dec-w 2.0 (2.92\,M, 0.883\,m). Interpolating the no-skip
      curve to equal parameter counts, skips are better by $0.031$\,m (at 2.09\,M) and
      $0.053$\,m (at 2.41\,M).

\item \textbf{Multi-scale features help further.} Adding an ASPP bottleneck to the skip model
      gives the best result of all (mean $0.791$\,m, P90 $1.432$\,m) --- another spatial, not
      capacity-driven, gain.

\item \textbf{Widening the encoder/bottleneck does not help.} The widest model (encoder-width
      1.4, 4.31\,M --- $79\%$ more parameters than the 2.41\,M skip model) is \emph{no better}
      (0.855 vs 0.830\,m). This is consistent with (2), though it is also confounded by being
      trained from scratch, so we treat it as supporting rather than decisive.
\end{enumerate}

% =====================================================================
\section{Does it help the thesis? --- Yes}
% =====================================================================
\begin{enumerate}[leftmargin=1.5em,itemsep=4pt]
\item \textbf{It justifies the chosen architecture.} The thesis selects the MobileNetV2-UNet as
      the final model because skip connections improve accuracy. Study A turns that choice from
      an empirical pick into a \emph{controlled, statistically significant} claim: skips are the
      active ingredient, and the gain is not an artifact of the extra parameters they add.

\item \textbf{It hardens the Mamba / 2-D spatial argument.} The thesis explains Mamba's failure
      (10.56\,m) as an \emph{architectural mismatch}: flattening the 2-D heatmap into a 1-D
      sequence destroys spatial locality. Study A provides the positive counterpart --- when
      spatial structure is \emph{preserved and enriched} (skips, then multi-scale), accuracy
      improves, while merely adding capacity does nothing. The failure of Mamba is therefore a
      \emph{spatial-information} problem, not a capacity problem --- exactly the thesis's claim,
      now backed by a parameter-controlled ablation.

\item \textbf{It points to a concrete future direction.} ASPP (multi-scale) gives the best
      result here, suggesting richer multi-scale spatial aggregation as the next architectural
      step --- a well-supported ``future work'' item.
\end{enumerate}

In short, Study A supplies the ``\emph{why}'' behind the headline result: compressed CSI
localization is driven by spatial-structure preservation, not model size.

% =====================================================================
\section{Caveats (for honesty in the write-up)}
% =====================================================================
\begin{itemize}[nosep,leftmargin=1.4em]
\item \textbf{Absolute level.} These means (e.g.\ $0.830$\,m for the skip dec-w 1.0 model) are a
      touch higher than the headline paper's $0.640$\,m \emph{median}, because (i) selection is
      by \emph{mean} (the median is pixel-quantized and cannot resolve the sub-decimetre gaps
      this study measures) and (ii) this is the configurable re-implementation. The study's
      value is in the \emph{relative} ladder comparisons, which are internally consistent.
\item \textbf{Single-seed configs.} Only the two decoder-width-1.0 anchors are multi-seed; the
      other seven configurations are single runs (direction/trend, not error bars). The
      decoder-width axis trend (flat no-skip, lower skip) is consistent across its points, and
      the one multi-seed pair is tight, which supports the conclusion.
\item \textbf{From-scratch secondary axis.} The encoder-width points (0.5, 1.4) lack ImageNet
      pretraining and are reported only as a caveated secondary axis; the clean capacity axis is
      decoder-width, where all points are pretrained.
\end{itemize}

% =====================================================================
\section{Conclusion}
% =====================================================================
On a controlled nine-architecture ladder, \textbf{spatial preservation --- not parameter count ---
drives WiFi-localization accuracy}. Adding decoder capacity to the plain MobileNetV2 leaves the
error flat ($\sim\!0.88$\,m across a $1.4\times$ parameter range), while skip connections lower it
at every capacity and even outperform larger no-skip models; multi-scale (ASPP) lowers it
further to $0.791$\,m. The skip-vs-no-skip gain is statistically significant
($p=1.8\times10^{-11}$) and seed-stable. This is exactly the result the supervisor asked for:
\emph{spatial preservation contributes independently of parameter count}. It validates the
thesis's architecture choice, converts the Mamba failure into a principled spatial-information
argument, and motivates multi-scale spatial aggregation as future work.

\end{document}
